\documentclass[11pt]{article}

\usepackage[a4paper, margin=1in]{geometry}
\usepackage{xcolor}
\usepackage{booktabs}
\usepackage{longtable}
\usepackage{hyperref}
\usepackage{verbatim}


\title{TAFS \,--\, Meezan Bank Bill Collection API \\ Integration Audit \& Reference Document}
\author{TAFS Development Team}
\date{8 June 2026}

\begin{document}
\maketitle

\section*{Purpose of This Document}
This document supersedes the previously shared integration guide and is prepared in direct
response to the points raised by the Meezan Bank IT Team. It provides:
\begin{itemize}
    \item Full request and response payload samples for \textbf{Bill Inquiry} and \textbf{Bill Payment}, covering both \textbf{success} and \textbf{failure} scenarios.
    \item A clear explanation that the fields returned for an \textbf{unsuccessful Bill Inquiry} are \textbf{structurally different} from those returned for a \textbf{successful Bill Inquiry} (different field names entirely \,--\, not just different values).
    \item A consolidated table of all \textbf{API response codes} used across both services.
    \item Direct clarifications against each point raised by the IT team.
\end{itemize}
This audit is based directly on the production source code of the integration
(\texttt{meezan.service.ts} / \texttt{meezan.controller.ts}), so every payload shown here is
guaranteed to match what the live system sends and receives.

\section{Base URL \& Authentication}

\subsection{Base URL}
\begin{quote}
\begin{verbatim}
https://tafs-backend-production.up.railway.app/api/v1/meezan
\end{verbatim}
\end{quote}

\subsection{Authentication Fields}
Every request (both Bill Inquiry and Bill Payment) must carry the following three
authentication fields in the JSON body. These are validated on \emph{every} call before any
business logic executes:

\begin{center}
\begin{tabular}{@{}lll@{}}
\toprule
\textbf{Field} & \textbf{Type} & \textbf{Description} \\
\midrule
\texttt{ServiceUserId} & string & Service authentication ID issued to Meezan \\
\texttt{UserPassword} & string & Password associated with the service user ID \\
\texttt{BillCompanyCode} & string & Company / campus identifier (\texttt{TAFCS}) \\
\bottomrule
\end{tabular}
\end{center}

\noindent\textbf{Validation order:}
\begin{enumerate}
    \item \texttt{ServiceUserId} and \texttt{UserPassword} are checked first. If either is wrong, the request is rejected with \texttt{ResponseCode "094"}.
    \item If \texttt{BillCompanyCode} is supplied and does not match, the request is rejected with \texttt{ResponseCode "095"}.
    \item Only after both checks pass does the system proceed to look up the voucher.
\end{enumerate}

\section{Service 1 \,--\, Bill Inquiry}
\textbf{Endpoint:} \texttt{POST /bill-inquiry} \hfill \textbf{HTTP Status:} Always \texttt{200 OK}

\noindent The HTTP status code is always \texttt{200}; the actual outcome (success or failure)
is communicated through the JSON body \,--\, see Section~2.4 for why this matters.

\subsection{Request Payload}
\begin{quote}
\begin{verbatim}
{
  "ServiceUserId": "0b986d818b54b6eeedbf7658",
  "UserPassword": "e67f36b15bd1d034017354db2486a82b",
  "BillCompanyCode": "TAFCS",
  "VoucherNumber": "05260003357"
}
\end{verbatim}
\end{quote}

\begin{center}
\begin{tabular}{@{}llll@{}}
\toprule
\textbf{Field} & \textbf{Type} & \textbf{Required} & \textbf{Description} \\
\midrule
\texttt{ServiceUserId} & string & Yes & Authentication ID \\
\texttt{UserPassword} & string & Yes & Authentication password \\
\texttt{BillCompanyCode} & string & Yes & Company / campus identifier \\
\texttt{VoucherNumber} & string & Yes & Voucher number to look up \\
\bottomrule
\end{tabular}
\end{center}

\subsection{Success Response \,--\, Voucher Found and Unpaid}
\textbf{HTTP 200}

\begin{quote}
\begin{verbatim}
{
  "StatusCode": "00",
  "StatusDesc": "Unpaid",
  "student_name": "TEST STUDENT B",
  "student_id": "2",
  "due_date": "20260521",
  "Amount_WID_Date": "50000",
  "Amount_AD_Date": "51000",
  "BillingMonth": "2504",
  "Remarks": "Late payment surcharge applies"
}
\end{verbatim}
\end{quote}

\begin{center}
\begin{tabular}{@{}ll@{}}
\toprule
\textbf{Field} & \textbf{Description} \\
\midrule
\texttt{StatusCode} & \texttt{"00"} \,--\, indicates a successful lookup \\
\texttt{StatusDesc} & \texttt{"Unpaid"} \,--\, the voucher exists and is payable \\
\texttt{student\_name} & Full name of the student \\
\texttt{student\_id} & Student CC (registration) number \\
\texttt{due\_date} & Due date, format \texttt{YYYYMMDD} \\
\texttt{Amount\_WID\_Date} & Amount payable \emph{within} (on or before) due date, in PKR \\
\texttt{Amount\_AD\_Date} & Amount payable \emph{after} due date (includes late surcharge), in PKR \\
\texttt{BillingMonth} & Billing period, format \texttt{YYMM} \\
\texttt{Remarks} & \texttt{"Surcharge waived"}, \texttt{"Late payment surcharge applies"}, or \texttt{""} \\
\bottomrule
\end{tabular}
\end{center}

\noindent\textit{Note:} \texttt{Remarks} is set to \texttt{"Surcharge waived"} when the voucher
has been manually exempted from late fees, \texttt{"Late payment surcharge applies"} when the
inquiry is being made after the due date, and an empty string \texttt{""} otherwise.

\subsection{Failure Response \,--\, Voucher Not Found / Invalid / Void / Paid / Expired / Auth Error}
\textbf{HTTP 200} (the transport-level status remains \texttt{200}; failure is indicated purely
by the body)

\begin{quote}
\begin{verbatim}
{
  "ResponseCode": "091",
  "ResponseDesc": "Voucher Id is invalid"
}
\end{verbatim}
\end{quote}

\begin{center}
\begin{tabular}{@{}ll@{}}
\toprule
\textbf{Field} & \textbf{Description} \\
\midrule
\texttt{ResponseCode} & 3-digit numeric code identifying the failure reason (see Section~4) \\
\texttt{ResponseDesc} & Human-readable description of the failure \\
\bottomrule
\end{tabular}
\end{center}

\subsection{IMPORTANT \,--\, Field Differences Between Success and Failure Responses}
\label{sec:field-diff}
This is the point specifically raised by the IT team and must be handled carefully on the
bank's parsing side: \textbf{a successful Bill Inquiry response and a failed/erroneous Bill
Inquiry response do not share the same field names.} They are two structurally distinct
payload shapes returned from the same endpoint:

\begin{center}
\begin{tabular}{@{}p{0.46\textwidth}p{0.46\textwidth}@{}}
\toprule
\textbf{Successful Inquiry \,--\, fields present} & \textbf{Failed Inquiry \,--\, fields present} \\
\midrule
\texttt{StatusCode} & \texttt{ResponseCode} \\
\texttt{StatusDesc} & \texttt{ResponseDesc} \\
\texttt{student\_name} & \emph{(none of the voucher/student fields are present)} \\
\texttt{student\_id} & \\
\texttt{due\_date} & \\
\texttt{Amount\_WID\_Date} & \\
\texttt{Amount\_AD\_Date} & \\
\texttt{BillingMonth} & \\
\texttt{Remarks} & \\
\bottomrule
\end{tabular}
\end{center}

\noindent\textbf{Practical guidance for the bank's integration layer:}
\begin{itemize}
    \item Do \emph{not} assume the voucher/student fields (\texttt{student\_name}, \texttt{due\_date}, etc.) will always be present \,--\, they are \emph{only} present when \texttt{StatusCode == "00"}.
    \item The presence of the \texttt{ResponseCode} key (rather than \texttt{StatusCode}) is itself the signal that the inquiry failed.
    \item Recommended parsing logic: check for \texttt{StatusCode} first \,--\, if present and equal to \texttt{"00"}, treat as success and read the voucher fields; otherwise, read \texttt{ResponseCode}/\texttt{ResponseDesc} and treat as a failure using the table in Section~4.
\end{itemize}

\section{Service 2 \,--\, Bill Payment}
\textbf{Endpoint:} \texttt{POST /bill-payment} \hfill \textbf{HTTP Status:} Always \texttt{200 OK}

\subsection{Request Payload}
\begin{quote}
\begin{verbatim}
{
  "ServiceUserId": "0b986d818b54b6eeedbf7658",
  "UserPassword": "e67f36b15bd1d034017354db2486a82b",
  "BillCompanyCode": "TAFCS",
  "VoucherNumber": "05260003357",
  "TransDate": "20260512",
  "TransAuthenticationCode": "AUTH123456",
  "TransAmount": "50000",
  "PaymentMode": "ONLINE",
  "BankCode": "MEZN",
  "BankName": "Meezan Bank Limited",
  "ChequeNo": "",
  "Status": "C",
  "ReasonCode": "",
  "DateOfReturn": "",
  "ReasonDescription": ""
}
\end{verbatim}
\end{quote}

\begin{center}
\begin{tabular}{@{}llll@{}}
\toprule
\textbf{Field} & \textbf{Type} & \textbf{Required} & \textbf{Description} \\
\midrule
\texttt{ServiceUserId} & string & Yes & Authentication ID \\
\texttt{UserPassword} & string & Yes & Authentication password \\
\texttt{BillCompanyCode} & string & No\textsuperscript{*} & Company / campus identifier \\
\texttt{VoucherNumber} & string & Yes & Voucher number being paid \\
\texttt{TransDate} & string & Yes & Transaction date, format \texttt{YYYYMMDD} \\
\texttt{TransAuthenticationCode} & string & Yes & Bank transaction reference / auth code \\
\texttt{TransAmount} & string & Yes & Amount paid, in PKR \\
\texttt{PaymentMode} & string & Yes & Payment mode (e.g. \texttt{CASH}, \texttt{CHQ}, \texttt{ONLINE}) \\
\texttt{BankCode} & string & Yes & Bank's branch / routing code \\
\texttt{BankName} & string & Yes & Name of the bank / branch \\
\texttt{ChequeNo} & string & No & Cheque number (cheque payments only) \\
\texttt{Status} & string & Yes & \texttt{C} = Cleared, \texttt{L} = Lodged, \texttt{R} = Returned \\
\texttt{ReasonCode} & string & No & Reason code (used when \texttt{Status = R}) \\
\texttt{DateOfReturn} & string & No & Return date, format \texttt{YYYYMMDD} (used when \texttt{Status = R}) \\
\texttt{ReasonDescription} & string & No & Reason description (used when \texttt{Status = R}) \\
\bottomrule
\end{tabular}
\end{center}
\noindent\textsuperscript{*}\,\texttt{BillCompanyCode} is optional in the payment request body;
however, if it \emph{is} supplied, it is still validated against our records and a mismatch
returns \texttt{ResponseCode "095"}.

\subsection{Status Field \,--\, Behaviour}
\begin{center}
\begin{tabular}{@{}lll@{}}
\toprule
\textbf{Status} & \textbf{Meaning} & \textbf{System action} \\
\midrule
\texttt{C} & Cleared \,--\, payment confirmed & Voucher marked \texttt{PAID}; deposit \\
 & & created; fee heads and (if applicable) \\
 & & overdue surcharges allocated and settled \\
\texttt{L} & Lodged \,--\, payment pending & No ledger update; logged only \\
\texttt{R} & Returned \,--\, payment rejected & No ledger update; reason logged only \\
\bottomrule
\end{tabular}
\end{center}
\noindent Only \texttt{Status = "C"} triggers actual ledger posting (deposit creation, fee
allocation, surcharge settlement, and marking the voucher \texttt{PAID}). \texttt{L} and
\texttt{R} are acknowledged with a \texttt{StatusCode "00"} response but make \emph{no} change
to the voucher or student ledger.

\subsection{Success Response \,--\, Payment Posted (\texttt{Status = "C"})}
\textbf{HTTP 200}
\begin{quote}
\begin{verbatim}
{
  "StatusCode": "00",
  "StatusDesc": "Success"
}
\end{verbatim}
\end{quote}
This confirms the voucher has been marked \texttt{PAID}, a deposit record has been created,
and all applicable fee heads (and overdue surcharges, where relevant) have been settled.

\subsection{Acknowledged but Not Posted \,--\, Lodged / Returned (\texttt{Status = "L"} or \texttt{"R"})}
\textbf{HTTP 200}
\begin{quote}
\begin{verbatim}
{
  "StatusCode": "00",
  "StatusDesc": "Lodged/Returned \xe2\x80\x94 not posted"
}
\end{verbatim}
\end{quote}
\noindent This is \emph{not} an error \,--\, \texttt{StatusCode "00"} confirms the request was
received and acknowledged. It simply means no ledger posting occurred because the
transaction was not in a \texttt{Cleared} state. The reason (for \texttt{Status = "R"}, including
\texttt{ReasonCode}/\texttt{ReasonDescription}) is recorded in our internal logs for audit
purposes.

\subsection{Failure Response \,--\, Voucher Not Found / Already Paid / Expired / Auth Error / Exception}
\textbf{HTTP 200}
\begin{quote}
\begin{verbatim}
{
  "ResponseCode": "097",
  "ResponseDesc": "Voucher already paid"
}
\end{verbatim}
\end{quote}
\noindent As with Bill Inquiry, failure responses use the \texttt{ResponseCode}/\texttt{ResponseDesc}
field pair rather than \texttt{StatusCode}/\texttt{StatusDesc}, and carry no payment/voucher
detail fields.

\section{Consolidated API Response Code Reference}
\label{sec:codes}
The following codes are returned (in the \texttt{ResponseCode} field, for failure cases) or as
\texttt{StatusCode "00"} for success/acknowledgement cases, across \textbf{both} Bill Inquiry
and Bill Payment:

\begin{center}
\begin{longtable}{@{}lp{0.30\textwidth}p{0.40\textwidth}@{}}
\toprule
\textbf{Code} & \textbf{Description} & \textbf{Returned by / Meaning} \\
\midrule
\endhead
\texttt{00} (\texttt{StatusCode}) & Success / Unpaid / Acknowledged & Both services. For Inquiry: voucher found and unpaid. For Payment: payment posted (\texttt{Status=C}), or acknowledged without posting (\texttt{Status=L/R}). \\[4pt]
\texttt{091} & Voucher Id is invalid & Both services. No voucher exists for the supplied \texttt{VoucherNumber} (including the 11-digit fallback lookup). \\[4pt]
\texttt{092} & Voucher date is expired & Both services. The voucher's \texttt{validity\_date} has passed relative to the inquiry/transaction date. \\[4pt]
\texttt{093} & Voucher is void & Bill Inquiry only. The voucher has been administratively voided. \\[4pt]
\texttt{094} & User invalid & Both services. \texttt{ServiceUserId} or \texttt{UserPassword} did not match our records. \\[4pt]
\texttt{095} & Company code mismatch & Both services. \texttt{BillCompanyCode} was supplied but did not match our records. \\[4pt]
\texttt{096} & General exception & Both services. An unexpected server-side error occurred while processing the request. \\[4pt]
\texttt{097} & Voucher is already paid / Voucher already paid & Both services. The voucher's status is already \texttt{PAID}. \\
\bottomrule
\end{longtable}
\end{center}

\noindent\textbf{Notes on code usage:}
\begin{itemize}
    \item Codes \texttt{091}, \texttt{092}, \texttt{094}, \texttt{095}, \texttt{096}, and \texttt{097} are common to both Bill Inquiry and Bill Payment.
    \item Code \texttt{093} (\textit{Voucher is void}) is returned \textbf{only} by Bill Inquiry. Bill Payment does not currently perform a separate void check (a voided voucher would typically already be reflected through one of the other statuses).
    \item All response codes are returned with \textbf{HTTP 200} \,--\, the bank's integration layer must inspect the JSON body (\texttt{StatusCode} vs. \texttt{ResponseCode}) to determine the outcome, not the HTTP transport status.
\end{itemize}

\section{Voucher Number Resolution Logic}
For both services, the voucher is first looked up by an exact match on \texttt{voucher\_number}.
If no match is found \emph{and} the supplied \texttt{VoucherNumber} is exactly 11 characters
long, the system attempts a fallback lookup by treating the last 7 characters as a numeric
internal voucher ID (\texttt{VoucherNumber.slice(4)}). This allows older/short-format voucher
numbers issued by the bank to still resolve correctly. If neither lookup succeeds,
\texttt{ResponseCode "091"} (\textit{Voucher Id is invalid}) is returned.

\section{Overdue Surcharge Logic}
\begin{itemize}
    \item \texttt{Amount\_WID\_Date} \,--\, payable amount if settled \emph{on or before} the due date.
    \item \texttt{Amount\_AD\_Date} \,--\, payable amount if settled \emph{after} the due date (includes the applicable late-payment surcharge).
    \item Whether a transaction is treated as overdue is determined by comparing the
    transaction/inquiry date to the voucher's \texttt{due\_date} (both normalized to midnight,
    so same-day payments are \emph{not} considered overdue).
    \item If a voucher has been manually marked as \texttt{surcharge\_waived}, no surcharge is
    applied regardless of the date, and \texttt{Remarks} will read \texttt{"Surcharge waived"}.
    \item During Bill Payment with \texttt{Status = "C"} on an overdue transaction, any
    un-waived arrear surcharges linked to the voucher are automatically settled alongside the
    regular fee heads in the same transaction.
\end{itemize}

\section{Clarifications on Points Raised by the IT Team}
\begin{enumerate}
    \item \textbf{``The previous document does not include API request and response samples.''} \\
    Resolved \,--\, Sections~2.1, 2.2, 2.3, 3.1, 3.3, 3.4 of this document now provide complete,
    code-accurate JSON samples for every request and response variant (success and failure)
    of both Bill Inquiry and Bill Payment.

    \item \textbf{``The API response codes are not mentioned.''} \\
    Resolved \,--\, Section~4 provides a single consolidated table of every response code
    (\texttt{00}, \texttt{091}--\texttt{097}) returned by either service, with its meaning and
    which service(s) return it.

    \item \textbf{``The fields received for an unsuccessful Inquiry are different from those in
    a successful Inquiry.''} \\
    Confirmed and documented in detail in Section~2.4. This is by design: a successful
    inquiry returns voucher/student details under \texttt{StatusCode}/\texttt{StatusDesc} plus
    the amount and billing fields, whereas a failed inquiry returns only
    \texttt{ResponseCode}/\texttt{ResponseDesc} with no voucher data. We recommend the bank's
    parser branch on the presence of \texttt{StatusCode == "00"} to distinguish the two shapes,
    as detailed in Section~2.4.
\end{enumerate}

\section{Recommended Integration Flow}
\begin{enumerate}
    \item \texttt{POST /bill-inquiry} with \texttt{VoucherNumber}.
    \item Check whether the response contains \texttt{StatusCode == "00"}:
    \begin{itemize}
        \item If yes \,--\, voucher is valid and unpaid; compare today's date against
        \texttt{due\_date} to decide whether to use \texttt{Amount\_WID\_Date} or
        \texttt{Amount\_AD\_Date}.
        \item If no \,--\, read \texttt{ResponseCode}/\texttt{ResponseDesc} and handle per
        Section~4 (e.g. do not proceed to payment if the voucher is invalid, void, expired,
        or already paid).
    \end{itemize}
    \item \texttt{POST /bill-payment} with \texttt{VoucherNumber}, the correct
    \texttt{TransAmount}, and \texttt{Status = "C"}.
    \item Check the response: \texttt{StatusCode == "00"} with \texttt{StatusDesc == "Success"}
    confirms posting; \texttt{StatusDesc == "Lodged/Returned \,--\, not posted"} confirms
    acknowledgement without posting; a \texttt{ResponseCode} indicates a failure to be handled
    per Section~4.
\end{enumerate}

\end{document}
